\subsection{Expressions}

\subsubsection{Value Categories}

\begin{tabularx}{\linewidth}{lX}
    \texttt{lvalue}     & Has name. Can take address.\\
    \texttt{prvalue}    & Temporary value. No name. \\
                        & \textcolor{magenta}{Guaranteed copy elision.} \\
    \texttt{xvalue}     & Expiring. Eg: \texttt{std::move(x)}.\\
    \texttt{glvalue}    & \texttt{lvalue} $\cup$ \texttt{xvalue}.\\
    \texttt{rvalue}     & \texttt{prvalue} $\cup$ \texttt{xvalue}.\\
\end{tabularx}

\begin{tabularx}{\linewidth}{lX}
    \texttt{L-value reference} & Refers to an \texttt{lvalue}. Eg: \texttt{int\& ref = var;}\\
    \texttt{R-value reference} & Refers to an \texttt{rvalue}. Eg: \texttt{int\&\& rref = 5;}\\
\end{tabularx}

\begin{mdframed}[backgroundcolor=gray!10,linecolor=Firebrick4]
\textbf{Exceptions:}\\
\vspace{-1.5em}
\begin{itemize}
\setlength\itemsep{-0.25em}
    \item R-value can bind to const l-value reference.\\
            Eg: \texttt{const int\& x = 42;}
    \item String literals are l-values.
\end{itemize}
\end{mdframed}

\subsubsection{Literals}

\begin{tabularx}{\linewidth}{llX}
    \texttt{Boolean}    & \texttt{true}                  & \\
                        & \texttt{false}                 & \\
    \texttt{Integer}    & \texttt{42}                    & Decimal. \\ 
                        & \texttt{4'200'000}             & Decimal with separators.\\
                        & \texttt{0b101010}              & Binary.\\
                        & \texttt{052}                   & Octal.\\
                        & \texttt{0x2a}                  & Hexadecimal.\\ 
    \texttt{Float}      & \texttt{3.14}                  & Double.\\
                        & \texttt{3.14f}                 & Float.\\ 
                        & \texttt{6.626e-34}             & \texttt{xey} $\to x\times 10^y$\\
    \texttt{Character}  & \texttt{'a'}                   & Ordinary literal.\\
                        & \texttt{'\textbackslash n'}    & Using escape seq.\\
                        & \texttt{'\textbackslash 13'}   & \\
                        & \texttt{u8'a'}                 & UTF-8.\\
    \texttt{String}     & \texttt{"Hello"}               & String literal.\\ 
    \texttt{Pointer}    & \texttt{nullptr}.              & \\   
\end{tabularx}
NOTE: \texttt{Literals are prvalues.}\\
\textcolor{magenta}{Exception: \texttt{String literals are lvalues.}}\\

\subsubsection{\texttt{std::initializer\_list}}
Constructed automatically from braced-init-list.\\
Eg:
\begin{mdframed}[nobreak=true,backgroundcolor=gray!10,linecolor=Firebrick4]
\begin{verbatim}
Foo foo({1, 2, 3});
print_list({1, 2, 3});
\end{verbatim}
\end{mdframed}

Passed as argument to functions.\\
Eg:
\begin{mdframed}[nobreak=true,backgroundcolor=gray!10,linecolor=Firebrick4]
\begin{verbatim}
void print\_list(std::initializer_list<int> list); 
Foo(std::initializer_list<int> list);
\end{verbatim}
\end{mdframed}

Access elements using iterators or range-based for loop.\\


\vfill\null
\columnbreak

\subsubsection{Arithemetic and logical operation}

    \begin{tabularx}{\linewidth}{l|lX}
        \hline\\
        & Operators & Associativity \\
        \hline \\
        1 & \texttt{() [] -$>$} \textbf{.} & left to right\\
        2 & \texttt{! $\tilde{}$ ++ -- + - * \& (type) sizeof} & right to left\\
        3 & \texttt{* / \%}(reminder) & left to right\\
        4 & \texttt{+ -} & left to right \\
        5 & \texttt{$<<$ $>>$} & left to right \\
        6 & \texttt{$<$ $<=$ $>$ $>=$}  & left to right\\
        7 & \texttt{== !=} & left to right\\
        8 & \texttt{\&} & left to right \\
        9 & \texttt{\^} & left to right \\
        10 & \texttt{|} & left to right\\
        11 & \texttt{\&\&} & left to right\\
        12 & \texttt{||} & left to right \\
        13 & \texttt{?:} & left to right \\
        14 & \texttt{= += -= *= /= \%= \&= \^{}= |= $<<$= $>>$=} & right to left\\
        15 & \texttt{,} & left to right\\
        \hline
    \end{tabularx}


    \textbf{Selected operators}\\
    \begin{tabularx}{\linewidth}{lX}
        \texttt{>>}             & Bitwise right shift operator.\\
                                & \texttt{a = b >> 1} equivalent to dividing \texttt{b} by 2.\\
        \texttt{<<}             & Bitwise left shift operator.\\
                                & \texttt{a = b << 1} equivalent to multiplying \texttt{b} by 2.\\
        \texttt{?:}             & Ternary conditional operator.\\
                                & \texttt{x = (condition)? expr1: expr2;}\\
                                & \texttt{expr1} $\to$ \texttt{condition == true}.\\
                                & \texttt{expr2} $\to$ \texttt{condition == false}.\\
        \texttt{$\tilde{}$}     & Bitwise negation.\\
        \texttt{\^}             & Bitwise XOR.\\
    \end{tabularx}


\subsubsection{Type casting}

    \begin{tabularx}{\linewidth}{lX}
        \texttt{const\_cast}            & Modifies \texttt{const}/\texttt{volatile} qualifiers.\\
                                        & Modifying a \texttt{const} object leads to UB.\\
        \texttt{static\_cast}           & Compile-time downcast.\\
                                        & Not safe.\\
        \texttt{dynamic\_cast}          & Safe downcast.\\
                                        & Eg: \texttt{Base* b = static\_cast<Base*>d;} \\
                                        & when invalid \texttt{b} is \texttt{nullptr}.\\
                                        & Eg: \texttt{Base\& b = static\_cast<Base\&>d;} \\
                                        & when invalid throws \texttt{std::bad\_cast}.\\
        \texttt{reinterpret\_cast}      & \\
    \end{tabularx}

